\begin{table}[h!]
\centering
\begin{threeparttable}
\caption{Summary Statistics ($n = 48$)}
\label{tab:sumstats}
\begin{tabular}{lcccccc}
\toprule
\textbf{Variable} & \textbf{Mean} & \textbf{SD} & $p_{5}$ & $p_{95}$ & \textbf{Min} & \textbf{Max} \\
\midrule
\multicolumn{7}{l}{\textit{Demographics}} \\[2pt]
Age &  27.5 &   3.6 &    24 &    34 &    24 &    43 \\
Female &  0.38 &  0.49 &  0.00 &  1.00 &     0 &     1 \\
Married &  0.10 &  0.31 &  0.00 &  1.00 &     0 &     1 \\
Has children &  0.08 &  0.28 &  0.00 &  1.00 &     0 &     1 \\[4pt]
\multicolumn{7}{l}{\textit{Wages (monthly GHC)}} \\[2pt]
Minimum WTA &   2248 &   1624 &    100 &   4000 &     18 &  10000 \\
Maximum WTA &   4260 &   2647 &    250 &  10000 &     30 &  10000 \\
\bottomrule
\end{tabular}
\begin{tablenotes}[flushleft]
\scriptsize
\item \textit{Notes:} Female, Married, and Has children are binary (0/1) indicators; their means equal the share of the sample with that characteristic. $p_5$ and $p_{95}$ for wage variables are the winsorization thresholds used in Table~\ref{tab:regsw}. Wages standardized to monthly GHC equivalents using a 40-hour, 5-day working week.
\end{tablenotes}
\end{threeparttable}
\end{table}
